\documentclass[11pt,a4paper]{article}
\usepackage[margin=2.5cm]{geometry}
\usepackage{amsmath,amssymb}
\usepackage{booktabs}
\usepackage{enumitem}
\usepackage[hidelinks]{hyperref}
\setlist{itemsep=2pt,topsep=4pt}
\newcommand{\code}[1]{\texttt{#1}}

\title{Bimodal experiment: MGD vs.\ regularised estimator\\[2pt]\large Recap of the changes of 24 September 2026}
\author{Chiara Zelco}
\date{24 September 2026}

\begin{document}
\maketitle

\section*{Summary}
The bimodal experiment compares two ways of estimating the parameters $\theta$ of a
maximum-entropy model: raw MGD, and a \emph{regularised} (REG) version which smooths the
estimate $\theta(t)$ along the SDE time. For each estimator we measure how much the fitted
energy fluctuates from one independent run to another, and compare it to the
Cram\'er--Rao bound. The results obtained in June 2026 have to be redone, because of
problems in how the regularised estimate was computed. This note explains what was wrong,
what was changed, and where the new run stands. \textbf{No new scientific result is
reported here}: the full rerun was launched at the end of the day.

\section{The experiment in a few words}
\begin{itemize}
  \item Data: $n_1 = 10^6$ samples from the one-dimensional bimodal density
        $p(x) \propto \exp(-\beta x^4 + 5\beta x^2 + \tfrac{\beta}{2} x)$,
        for ten values $\beta = 0.1, 0.3, \dots, 1.9$.
  \item Model: four potentials $\phi_a(x) = x^a/a$, $a = 1,\dots,4$. The true parameters
        $\theta^*$ are therefore known exactly, which makes this a good test case.
  \item For each $\beta$ we repeat the whole MGD procedure $K = 100$ times with different
        random seeds (new data and new initial noise each time), with $n_t = 50\,000$ SDE steps and $\sigma = 10$.
  \item From the $K$ final estimates $\hat\theta$ we compute the variance and the error of
        the fitted energy, for MGD and for REG, and compare with the Cram\'er--Rao bound.
\end{itemize}

\section{What was wrong in the June runs}
\begin{enumerate}
  \item \textbf{Time averaging with a shift.} All June runs grouped the SDE steps in blocks of
        10 (\code{n\_subsample = 10}) before regularising. Each block was labelled with the time
        of its \emph{first} step although it contains the average of ten steps, so the
        regularised curve was shifted in time by half a block.
  \item \textbf{Low-precision solve.} The regularised problem was solved as one large linear system
        in single precision (float32). We had already seen elsewhere that this solve becomes
        unreliable when the matrices are badly conditioned and the smoothing is strong.
  \item \textbf{Smoothing strength fixed in advance.} The smoothing parameter $\lambda$ was chosen
        by hand before each run ($10^{-3}$ to $10^{-7}$ depending on the sweep), with no way of
        checking afterwards whether it was too strong or too weak.
  \item \textbf{Two conventions mixed in the variance.} The code estimates the coefficients of
        $x^a/a$, but the Cram\'er--Rao bound and the energy variance were computed with the
        potentials $x^a$. The variances of $\theta$ and their bound were thus off by a factor
        $a^2$ per component, and the energy variance mixed the two conventions. This is
        independent of the regularisation issue, and it affects the June variance plots.
  \item \textbf{A diverging term at $t = 1$.} One quantity of the regularised system contains
        $\tan(\pi t/2)$, which is infinite at the last time step. It was never actually used,
        but it is now set to zero explicitly.
\end{enumerate}

\section{What was changed}

\subsection*{Organisation of the code}
\begin{itemize}
  \item The experiment folder was moved out of \code{turbulence/} to the top level of the
        repository: \code{bimodal\_experiment/}.
  \item Its helper file \code{utils.py} was renamed \code{bimodal\_utils.py}, so that it cannot be
        confused with the shared \code{codes/utils.py}. A dependency on a file missing from the
        repository (\code{Mala.py}) no longer breaks the import. In the notebooks, only the import line changed.
\end{itemize}

\subsection*{The regularised estimate}
\begin{itemize}
  \item No more time blocks: the regularisation now uses every SDE step (\code{n\_subsample = 1}).
        If blocks are ever used again, each block is labelled at its mean time.
  \item The linear system is solved in double precision with the same solver used for the
        turbulence experiments (a block-tridiagonal ``Thomas'' solver), instead of a separate copy.
  \item The run no longer solves for a single $\lambda$. It saves the ingredients of the
        regularised problem, so that it can be solved afterwards for any $\lambda$, without
        repeating the (expensive) SDE simulation.
  \item To make this fast, a version of the solver was written that treats all runs and all
        values of $\lambda$ at once. It gives the same answer as the original solver (relative
        difference $7\times10^{-13}$ on test problems) and is about 65 times faster: roughly
        13 minutes per $\beta$ instead of 15 hours.
\end{itemize}

\subsection*{Choosing $\lambda$ after the run}
For a grid of 12 values ($\lambda = 0$ and $10^{-8}$ to $10^{-3}$) we solve the regularised
problem and compare the smoothed curve $\Theta_\lambda(t)$ with the unsmoothed one
$\Theta_0(t)$. As long as $\lambda$ only removes noise, the difference between the two looks
like noise: its averages over blocks of 50 steps are small. When $\lambda$ becomes too large,
the difference contains a slow systematic part (real structure is being erased), and the test
statistic $E[z^2]$ rises above 1. \textbf{We keep the largest $\lambda$ with $E[z^2] \le 1$ in
every time window.} This is the same rule as for turbulence, and it does not use the true
$\theta^*$.

Because $\theta^*$ is known in this experiment, we also record the $\lambda$ that would be best
in hindsight (the ``oracle'', smallest error with respect to $\theta^*$). It is only used
to check whether the rule above picks a good value; it cannot be used for turbulence.

\subsection*{How the energy error is measured}
The fitted energy is $E(x) = \hat\theta \cdot \phi(x)$. Adding a constant to the energy does
not change the probability distribution, so we measure only the variations of the energy
\emph{shape}: the energy is centred over $x$ before comparing runs. With $C$ the covariance of
$\hat\theta$ over the $K$ runs, the energy variance is
\[
  \operatorname{Var}(E) = \operatorname{tr}\bigl(C\,\operatorname{Cov}_x[\phi]\bigr),
  \qquad
  \text{Cram\'er--Rao bound} = \frac{r}{n_1} = \frac{4}{n_1}.
\]
Two versions are reported side by side:
\begin{itemize}
  \item \textbf{variance}: runs compared with their own average (spread between runs only);
  \item \textbf{error (MSE)}: runs compared with the true $\theta^*$ (spread plus bias).
\end{itemize}
This distinction matters for REG: smoothing always reduces the variance, but it can introduce
a bias. The Cram\'er--Rao bound only applies to unbiased estimators, so a REG variance below
the bound would be a sign of bias, not of a better estimator. Comparing the two numbers shows
how much of the error is bias.

\section{Running on Jean Zay}
\begin{itemize}
  \item \textbf{Account error.} The first submission was refused (``invalid account or
        partition''). The job settings were replaced by ones known to work.
  \item \textbf{Memory.} The next attempt (H100) ran out of memory during the first run:
        memory grew by about 2\,MB per SDE step. The cause could not be identified with certainty;
        the only change inside the SDE loop since June (a small copy to CPU memory at every step)
        was removed, and memory is now written in the log every 200 steps.
  \item \textbf{Runtime.} Done one after the other, the 100 runs of one $\beta$ would take
        about 11 hours (the June runs took 19\,h). The runs are now spread over
        100 parallel jobs on V100 GPUs (10 values of $\beta$ $\times$ 10 groups of 10 runs,
        about 2 hours each), followed by one short job per $\beta$ that collects the runs and
        chooses $\lambda$. Each run saves its own files, so failed jobs can be resubmitted
        without redoing finished runs, and splitting the runs gives exactly the same results
        as running them in sequence.
\end{itemize}

\section{Test run}
Before the full launch, a small test was run on Jean Zay ($\beta = 0.5$, $K = 2$ runs,
$n_t = 5000$ steps). It was meant to check that everything runs, \textbf{not} to produce results:
a variance from two runs means nothing.
\begin{itemize}
  \item Memory stayed flat (0.97\,GB at step 200, 0.98\,GB at step 4800): the leak is gone.
  \item Speed: about 13\,ms per step, i.e.\ about 11 minutes for a full run of 50\,000 steps.
  \item Checks: with $\lambda = 0$ the regularised estimate coincides with raw MGD, as it must;
        the Cram\'er--Rao value is exactly $4/n_1 = 4\times10^{-6}$; the linear solves are
        accurate to $10^{-11}$.
  \item The $\lambda$ chosen by the rule ($3\times10^{-5}$) happened to equal the oracle one.
        With two short runs this is encouraging but not evidence.
\end{itemize}

\section{Status and next steps}
\begin{enumerate}
  \item The full sweep ($10$ values of $\beta$, $K = 100$, $n_t = 50\,000$) was launched on
        24 September; results are expected a few hours later.
  \item Collect the results and produce the figures (energy variance and error vs.\ $\beta$ for
        MGD, REG with the chosen $\lambda$, REG with the oracle $\lambda$, and the bound; error vs.\
        $\lambda$ for each $\beta$).
  \item Check whether the rule picks a $\lambda$ close to the oracle one. If it does, this
        supports using the same rule for turbulence, where no true $\theta$ is available.
  \item Commit the changes (they are not yet in the git history).
\end{enumerate}

\appendix
\section*{Appendix: files}
{\small\begin{tabular}{@{}ll@{}}
\toprule
File & Role \\
\midrule
\code{bimodal\_experiment/bimodal\_ensemble\_reg.py} & runs MGD, saves each run; collects the runs \\
\code{bimodal\_experiment/sde\_routines\_scalar\_reg.py} & 1-D SDE; builds the regularised problem \\
\code{bimodal\_experiment/select\_lambda.py} & solves for all $\lambda$, chooses $\lambda$ \\
\code{bimodal\_experiment/launch/bimodal\_sweep.sh} & submits the jobs on Jean Zay \\
\code{bimodal\_experiment/analyse\_lamsweep.ipynb} & figures \\
\code{codes/lam\_selection.py} & fast solver and the $\lambda$ rule (shared) \\
\bottomrule
\end{tabular}}

\end{document}
